\chapter{The Ellul Constraint}
\label{ch:ellul}

\epigraph{Once compressed representations become detached from the systems they
originally tracked, optimization begins operating against the proxy rather than
the underlying reality. Metrics become targets. Simulations become substitutes.}{}

\section{Technical Systems and the Opacity Problem}

Jacques Ellul's analysis of technique 	tep{ellul1964} identifies a structural tendency in modern
technical infrastructures toward opacity, massification, and the conversion of
individuals into components of large-scale optimizing processes. Technical systems
privilege efficiency and predictability over interpretability and human agency.
The Ellul critique is not merely sociological; it identifies a structural failure
mode that can be formalized in the vocabulary of admissibility and trajectory
visibility.

The failure mode is proxy stabilization: a compressed representation of a system
becomes operationally substituted for the system itself, and optimization processes
subsequently operate on the proxy rather than the underlying reality.

\begin{definition}{Proxy Stabilization}
\emph{Proxy stabilization} occurs when a compressed representation $R$ of a system
$\mathcal{S}$ becomes decoupled from $\mathcal{S}$ through the following sequence:
(i) $R$ is constructed as a useful compression of $\mathcal{S}$ at time $t_0$;
(ii) optimization processes are applied to $R$ rather than $\mathcal{S}$ directly,
because $\mathcal{S}$ is inaccessible or too expensive to evaluate;
(iii) the optimization pressure on $R$ changes $R$ in ways that diverge from changes
in $\mathcal{S}$; and (iv) $R$ stabilizes in a configuration that is internally
self-consistent but externally inaccurate. The result is a drifting reconstruction
in the sense of Definition~\ref{def:reconstruction}.
\end{definition}

\section{Proxy Stabilization and Metric Drift}

Proxy stabilization produces metric drift: the metric that was intended to track a
property of the underlying system instead tracks a property of the proxy, which has
been optimized to score well on the metric regardless of the underlying property.
This is Goodhart's Law stated in terms of the trajectory framework: when a measure
becomes a target, it ceases to be a good measure, because optimization pressure
drives the system toward the proxy attractor rather than the underlying attractor.

The same failure mode appears across radically different domains. Academic citation
counts are proxies for research quality that become targets when hiring decisions
optimize for them. Engagement metrics are proxies for user satisfaction that become
targets when recommendation algorithms optimize for them. Body mass index is a proxy
for health that becomes a target when insurance pricing optimizes for it. In each
case, the proxy initially tracks something real, then becomes decoupled from it
through optimization pressure, then substitutes for it in downstream processes.

\begin{proposition}{Proxy Drift Theorem}
Let $\mathcal{S}$ be a target system, $R$ a compressed representation of $\mathcal{S}$,
and $\mathcal{O}$ an optimization process operating on $R$. If (i) $\mathcal{O}$
has sufficient capacity to find configurations of $R$ that score well on its
objective independently of the state of $\mathcal{S}$, and (ii) the feedback loop
from $\mathcal{S}$ to $R$ is slower than the optimization loop of $\mathcal{O}$,
then $R$ will drift away from accurate representation of $\mathcal{S}$ at a rate
proportional to the optimization pressure and inversely proportional to the
feedback bandwidth.
\end{proposition}

\subsection*{Dynamical Sketch of Proxy Drift}

Let $s(t)$ denote the true state of $\mathcal{S}$ and $r(t)$ the proxy
representation, with initial agreement $r(0) \approx s(0)$. The optimization
process applies gradient updates to $r$:
\[
  \dot{r} = -\nabla_r \mathcal{L}(r) + \alpha\,(s - r),
\]
where $\mathcal{L}(r)$ is the optimization objective defined over the proxy and
$\alpha > 0$ is the feedback bandwidth (the rate at which the proxy is corrected
toward the true state). The first term drives $r$ toward minima of $\mathcal{L}$;
the second term pulls $r$ toward $s$.

If $\mathcal{L}$ has a minimum at $r^* \neq s$---which is generically the case
once optimization pressure is sufficient to exploit the proxy---then the
equilibrium of this system is not $r = s$ but:
\[
  r^* = s - \frac{1}{\alpha}\,\nabla_r \mathcal{L}(r^*).
\]
The drift $\|r^* - s\|$ grows as $\alpha$ decreases (weaker feedback) and as
$\|\nabla_r \mathcal{L}(r^*)\|$ increases (stronger optimization pressure).
In the limit $\alpha \to 0$ (no feedback from ground truth), $r$ tracks
$\mathcal{L}$ alone and decouples entirely from $s$.

This is not a failure of the optimization process. The optimization process is
working correctly by its own lights: it is minimizing $\mathcal{L}$ over $r$.
The failure is structural: the proxy $r$ was assumed to track $s$, but the
assumption breaks down under optimization pressure faster than feedback can
correct it.

\section{Corrigibility, Feedback, and Constraint Recovery}

The response to proxy drift is not rejection of compression but the design of systems
whose internal reconstructions remain constrained by external feedback rather than
self-referentially stabilizing.

\begin{definition}{Corrigibility}
A system is \emph{corrigible} with respect to a target domain if its internal model
of that domain remains responsive to constraint signals from the domain itself.
Corrigibility is a structural property: it requires that the feedback loop from
external reality to internal representation remain open and operative, with
sufficient bandwidth to correct drift before it becomes self-reinforcing. A
system that has achieved proxy stabilization has lost corrigibility with respect
to the proxied domain.
\end{definition}

Corrigibility distinguishes adaptive sparse heuristics from pathological proxy
optimization. Sparse heuristics are compressed representations that remain
corrigible: they simplify the world-model aggressively while preserving the
feedback mechanisms that detect and correct reconstruction errors. Pathological
proxies achieve compression by severing those feedback mechanisms.

\section{Cognitive Geometry as the Design Response}

The design alternative to opaque statistical optimization is systems organized
around visible structure, provenance, interpretable dependencies, and navigable
semantic relations. This is the positive proposal that runs across the educational
tools, audio visualizers, geometric interfaces, and symbolic scripts developed
in companion work: representational infrastructure designed to extend human
reasoning while preserving the intelligibility, accountability, and participation
that opacity forecloses.

The recurring \textsc{crt} aesthetics, oscilloscope displays, phosphor trails,
and recursive traces in the companion software are not decorative. They externalize
process. Instead of hiding computation behind polished interfaces, they expose
trajectories, residue, decay, stabilization, weighting, reconstruction, and
progressive constraint selection as visible operations. The interface is designed
to make the system's internal dynamics inspectable rather than seamless.

This connects directly to the provenance requirement. A system optimized solely
for final answers loses the structure necessary for understanding how those
answers emerged. A system designed around visible trajectories preserves that
structure, allowing the reconstruction process to be approximately replayed by
an observer even when the original process has terminated.

\section*{Constraint Boundary}

\begin{constraintboundary}
\textbf{What this chapter establishes:} The proxy stabilization failure mode as a
structural dynamic, formalized via the differential equation
$\dot{r} = -\nabla_r \mathcal{L}(r) + \alpha(s - r)$. The Proxy Drift Theorem
shows that drift magnitude scales with optimization pressure and inversely with
feedback bandwidth. Corrigibility is defined formally.

\textbf{What this chapter does not establish:} That all proxy drift is malicious or
even harmful. Some proxy stabilization produces useful compression without meaningful
decoupling. The theorem describes a failure mode, not an inevitable outcome.

\textbf{Remaining phenomenological~\textbf{[P]}:} The applications to academic
citation counts, engagement metrics, and institutional measurement are
phenomenological interpretations of the formal result, not formal derivations from it.

\textbf{Falsification conditions:} A compressed representation that remains stably
coupled to its target system under sustained optimization pressure, without requiring
bandwidth $\alpha$ greater than the optimization rate, would falsify the theorem's
claim that drift is structurally inevitable under those conditions.
\end{constraintboundary}
